\section{Subsequences}
%%%%%%%%%%%
%%%%%%%%%%%

\begin{definition}[subsequence]
\mymark{*}
If $a_n$ is a sequence and $n_k$ is a strictly increasing sequence whose values
are natural numbers, then the sequence
$b_k = a_{n_k}$ is called a \emph{subsequence}%
\mymargin{subsequence}%
\index{subsequence} of $a_n$
(or also a \emph{sequence extracted} from $a_n$).
\end{definition}

Recalling that a sequence $a_n$ is nothing but a function
$\vec a\colon \NN \to \RR$, the sequence $n_k$ corresponds to a function
$\vec n\colon \NN \to \NN$ and the subsequence $a_{n_k}$ corresponds to the
composite function $\vec a \circ \vec n$.

Note that in the previous definition the variable $n$ represents
a dummy variable when we write the sequence $a_n$, but
it also represents the name of the fixed sequence $n_k$.
This overload
of meaning is intentional and, if used correctly, simplifies
notations, as the sequence $n_k$ is substituted for the
variable $n$, with the same name, in the sequence $a_n$.
The subsequence $a_{n_k}$ turns out to be a sequence in the variable $k$, not in
the variable $n$.

\begin{example}
Let $a_n = n^2$ be the sequence of perfect squares:
\begin{center}
\begin{tabular}{l|rrrrrrrrr}
$n$   & $0$ & $1$ & $2$ & $3$ & $4$  & $5$  & $6$  & \dots \\ \hline
$a_n$ & $0$ & $1$ & $4$ & $9$ & $16$ & $25$ & $36$ & \dots
\end{tabular}
\end{center}
Consider the sequence of even numbers $n_k = 2k$.
The corresponding subsequence of perfect squares
$b_k = a_{n_k}$
represents the sequence of squares of even numbers:
\begin{center}
\begin{tabular}{l|rrrrrrrrr}
$k$       & $0$ & $1$ & $2$ & $3$ & $4$  & $5$  & $6$  & \dots \\ \hline
$n_k$ & $0$ & $2$ & $4$ & $6$ & $8$ & $10$ & $12$ & \dots \\
$a_{n_k}$ & $0$ & $4$ & $16$ & $36$ & $64$ & $100$ & $144$ & \dots
\end{tabular}
\end{center}
In practice, we have
  $b_k = a_{n_k} = a_{2k} = (2k)^2$.

We have indeed \emph{extracted} some of the terms of the original sequence.
\end{example}

\begin{example}
If $a_n = (-1)^n$ and $n_k=2k$ then $a_{n_k} = 1$.
We see that an irregular sequence
can contain a regular subsequence.
\end{example}

Observe that if $\vec n\colon \NN\to \NN$ is a
strictly increasing function
(that is, $n_k=\vec n(k)$ is a strictly increasing sequence of indices)
then setting $A=\vec n(\NN)=\ENCLOSE{n_k\colon k\in \NN}$ we have that
$\vec n \colon \NN \to A$ is a bijection. Thus $A$ is an infinite set.
Conversely, given any infinite set $A\subset \NN$, there exists a
unique sequence $\vec n\colon \NN \to A$ that is bijective and strictly
increasing:
it suffices to set, by induction,
$n_0 = \min A$, $n_1 = \min \ENCLOSE{n\in A \colon n > n_0}$
and, in general,
 $n_{k+1} = \min\ENCLOSE{ n\in A \colon n > n_k}$.

Thus we can identify the subsequences of a sequence
$\vec a \colon \NN \to \RR$ with the restrictions to infinite subsets of $\NN$.
In the previous example, we considered the subsequence
of all terms with even indices $n_k=2k$ to obtain the subsequence
$a_{n_k} = a_{2k}$. We can equivalently think of taking the set of
all even numbers $A=2\NN$ and consider the sequence restricted to only
even indices:
\[
  a_0, a_2, a_4, \dots
\]
If we renumber the even indices using all natural numbers, we obtain
the subsequence $b_k=a_{n_k}$:
\[
  b_0 = a_0,\ b_1 = a_2,\ b_2 = a_4,\ \dots,\ b_k = a_{n_k},\ \dots
\]

\begin{lemma}[monotone subsequences]%
\label{lem:estratte_monotone}%
Every sequence $a_n\in \RR$ has a monotone subsequence $a_{n_k}$.
Moreover, if $\sup a_n = +\infty$, there is a strictly increasing subsequence
$a_{n_k}$
that tends to $+\infty$.
\end{lemma}
%
\begin{proof}
Consider the set $P$ of "peak" points, i.e., the indices
of those terms of the sequence that are greater than or equal to all subsequent
terms:
\[
  P = \ENCLOSE{n\in \NN\colon m\ge n \implies a_n\ge a_m}.
\]
If $P$ is finite,
it means that there exists an index $n_1\in \NN$ such
that there are no peaks from $n_1$ onwards. In particular, $n_1$ is not a peak
point,
so there must exist $n_2>n_1$ such that $a_{n_2}>a_{n_1}$.
But neither is $n_2$ a peak point, so there must exist $n_3>n_2$ such
that $a_{n_3}>a_{n_2}$... proceeding inductively, we can thus define
a sequence $n_k$ of indices such that $a_{n_k}$ turns out to be strictly
increasing.

In particular, if $\sup a_n = +\infty$, we are in the previous situation
because clearly, in this case, $P$ is empty since for every $n\in \NN$,
there must exist $m\in \NN$ such that $a_m > \max\ENCLOSE{a_0, a_1, \dots a_n}$
and certainly $m>n$.

If, on the other hand, $P$ is infinite, then listing its elements in order, we
will obtain
a sequence $n_1 < n_2 < n_3, \dots$ of indices each of which corresponds to a
peak value.
If $j>k$, then $n_j>n_k$, and since $n_k\in P$, it means that
$a_{n_j} \le a_{n_k}$.
Thus the sequence $a_{n_k}$ turns out to be decreasing.
\end{proof}


\begin{theorem}[Bolzano-Weierstrass]\label{th:Bolzano}
\label{th:bolzano_weierstrass}%
\mymark{***}%
\mymargin{Bolzano-Weierstrass}%
\index{Bolzano-Weierstrass}%
\index{theorem!of Bolzano-Weierstrass}%
Every sequence $a_n$ has a regular subsequence.
More precisely, if $a_n$ is a bounded sequence,
then there exists a convergent subsequence
$a_{n_k}$.
If $a_n$ is an unbounded sequence, then
there exists a divergent subsequence $a_{n_k}$.
\end{theorem}
%
\begin{proof}
\mymark{***}
Let us start with the real case $a_n\in \RR$.
Lemma~\ref{lem:estratte_monotone} guarantees
that there exists a monotone subsequence $a_{n_k}$.
But by Theorem~\ref{th:limite_successione_monotona}, we immediately conclude
that $a_{n_k}$ has a limit.
If $a_n$ is bounded, then $a_{n_k}$ is also bounded, and therefore
in this case, the limit is finite, and thus the extracted sequence is
convergent.

If $a_n$ is not bounded above, the lemma
guarantees that the extracted sequence tends to $+\infty$
and thus is divergent.
By symmetry (for example, considering $b_n=-a_n$), if $a_n$ is not bounded
below,
there exists an extracted sequence that tends to $-\infty$, and thus in
this case, the extracted sequence is divergent.

If $a_n$ is a sequence of complex numbers, we can write
$a_n = x_n + i y_n$ with
$x_n$ and $y_n$ being real sequences. If $a_n$ is bounded, it means that
$\abs{a_n}$
is bounded above. But it turns out that $\abs{x_n} \le \abs{a_n}$ and
$\abs{y_n}\le \abs{a_n}$, so if $a_n$ is bounded, both the real part
$x_n$ and the imaginary part $y_n$ are bounded sequences.
Then $x_n$ admits a convergent subsequence $x_{n_k}$.
But $y_{n_k}$ is also bounded, and therefore it also admits a
sub-subsequence $y_{n_{k_j}}$ that is convergent.
Thus the sub-subsequence $a_{n_{k_j}}$ is convergent.

If $a_n\in \CC$ is not bounded, by applying the theorem to its modulus
$\abs{a_n}$,
we find that there is an extracted sequence $a_{n_k}$ such that
$\abs{a_{n_k}}\to +\infty$.
But this means that $a_{n_k}\to \infty \in \bar \CC$.
\end{proof}

%%%%%%%%%%%%
%%%%%%%%%%%%
\subsection{limit points}
%%%%%%%%%%%%
%%%%%%%%%%%%

\begin{theorem}[bridge between function limits and sequence limits]%
\label{th:ponte}%
\mymark{***}%
Let $A \subset \RR$, $f\colon A \to \RR$, let $x_0$ be an accumulation point of
$A$, and let
$\ell \in [-\infty, +\infty]$.
The following two conditions are equivalent:
\begin{enumerate}
\item $\displaystyle \lim_{x\to x_0} f(x) = \ell$;
\item for every sequence $a_n\to x_0$ with $a_n\in A\setminus\ENCLOSE{x_0}$, it holds that
\[
\lim_{n\to+\infty} f(a_n) = \ell. 
\]
\end{enumerate}
\end{theorem}
%
\begin{proof}
\mymark{***}
If for $x\to x_0$ we have $f(x)\to \ell$ and if $a_n \to x_0$ with $a_n\in
A\setminus\ENCLOSE{x_0}$, the sequence $f(a_n)$ is nothing but the composition
of the function $f$ with the function $n\mapsto a_n$. We can then apply
the theorem on the limit of the composite function to obtain that $f(a_n)\to
\ell$.

Conversely, suppose we know that for every sequence $a_n\to x_0$, we have
$f(a_n)\to \ell$. We want to show that $f(x)\to \ell$. We do this by
contradiction: suppose there exists a neighborhood $U$ of $\ell$ such that for
any neighborhood $V$ of $x_0$, it is not the case that
$f((A\setminus\ENCLOSE{x_0})\cap V)\subset U$.
We can consider, for each $n\in \NN$, neighborhoods $V_n$ that are increasingly
smaller. For example, in the case $x_0 \in \RR$, we could choose $V_n =
(x_0-1/n, x_0+1/n)$; in the case $x_0 = +\infty$, we could choose $V_n =
(n,+\infty]$; and in the case $x_0=-\infty$, we could choose $V_n = [-\infty,
-n)$.
If, by contradiction, $f((A\setminus\ENCLOSE{x_0}\cap V_n))$ were not contained
in $U$,
it would mean that for each $n\in\NN$, there exists $a_n \in
(A\setminus\ENCLOSE{x_0})\cap V_n$ such that $f(a_n)\not \in U$. But then $a_n$
would be a sequence in
$A\setminus \ENCLOSE{x_0}$ with limit $x_0$
(since for every neighborhood of $x_0$, there exists an $N$ such that $V_N$ is
contained in that neighborhood, and for every $n>N$, we have $V_n\subset V_N$)
but $f(a_n)$ could not have limit $\ell$
(since it is outside the neighborhood $U$).
But this contradicts the hypothesis and thus concludes the proof of the theorem.
\end{proof}
\begin{proposition}[characteristic property of convergence]
  \label{prop:convergenza}
Let $f\colon A\subset \RR\to \RR$ and $x_0$ be an accumulation point of $A$.
Let $\ell\in \bar \RR$.
If for every $a_n\to x_0$, $a_n\in A\setminus\ENCLOSE{x_0}$,
there exists a subsequence $a_{n_k}$ such that $f(a_{n_k})\to \ell$, 
then $f(x)\to \ell$ as $x\to x_0$.
\end{proposition}
%
\begin{proof}
  If, by contradiction, $f(x) \not\to \ell$,
  there would exist a neighborhood $U\in \B_\ell$
  such that frequently $f(x) \not \in U$.
  This means that there exists a sequence
  $a_n\to x_0$ with $a_n\neq x_0$ such that $f(a_n)\not \in U$.
  But then from $f(a_n)$,
  it is not possible to extract a subsequence
  that has limit $\ell$, and this is contrary to the
  hypotheses.
\end{proof}

\begin{definition}[limit points, upper limit, lower limit]
  Let $f\colon A\subset \RR \to \RR$ be a function and $x_0\in \bar \RR$ 
  an accumulation point for $A$.
  A quantity $\ell\in \bar\RR$ is said to be
  a \emph{limit point}%
\mymargin{limit point}%
\index{limit!point} of $f(x)$ as $x\to x_0$ if 
  for every $U$ neighborhood of $\ell$, $f(x)\in U$ 
  frequently as $x\to x_0$.
  
  If we denote by $L\subset \bar\RR$ the set
  of limit points, we can define the \emph{upper limit} and the
  \emph{lower limit}
  \mymargin{upper/lower limit}%
\index{limit!upper/lower}
  respectively as
  \[
  \limsup_{x\to x_0} f(x) = \sup L, \qquad
  \liminf_{x\to x_0} f(x) = \inf L.
  \]
\end{definition}
  
\begin{proposition}[countable base of neighborhoods]%
  \label{prop:base_numerabile}%
    For every $\ell\in \bar \RR$, there exists a sequence $U_n$ of neighborhoods
  of $\ell$
  with the following properties:
  \begin{enumerate}
    \item $U_{n+1}\subset U_n$;
    \item for every $U$ neighborhood of $\ell$, there exists $n\in \NN$ such that $U_n\subset U$;
    \item if $a_n\in U_n$, then $a_n\to \ell$.
  \end{enumerate}
\end{proposition}
\begin{proof}
We can explicitly define the neighborhoods $U_n$ as follows:
\[
  U_n = 
  \begin{cases}
    \openinterval{\ell-\frac 1 n}{\ell+\frac 1 n} & \text{if $\ell\in \RR$}\\ 
    \opencloseinterval{n}{+\infty} & \text{if $\ell=+\infty$}\\
    \closeopeninterval{-\infty}{-n} & \text{if $\ell=-\infty$.}
  \end{cases}
\]
The first property is obviously satisfied.
The Archimedean property of real numbers guarantees the validity of the second 
property, from which the third follows immediately.
\end{proof}

\begin{proposition}[characterization of limit points]
  Let $f\colon A\subset \RR\to \RR$, $x_0\in \bar \RR$ 
  be an accumulation point for $A$. 
  The following are equivalent:
  \begin{enumerate}
    \item $\ell$ is a limit point of $f(x)$ as $x\to x_0$.
    \item There exists a sequence $a_n\in A$, $a_n\neq x_0$, $a_n\to x_0$ 
    such that $f(a_n)\to \ell$.
  \end{enumerate}
\end{proposition}
%
\begin{proof}
    Let $U_n$ be neighborhoods of $\ell$ and $V_n$ be neighborhoods of $x_0$
  defined
  as in Proposition~\ref{prop:base_numerabile}.
  If $\ell$ is a limit point of $f(x)$ as $x\to x_0$, 
  it means that for every $U$ neighborhood of $\ell$ and for every $V$ 
  neighborhood of $x_0$, there exists $x\in A\cap V\setminus\ENCLOSE{x_0}$ 
  such that $f(x)\in U$.
  Thus for every $n$, there exists $a_n\in A \cap V\setminus\ENCLOSE{x_0}$ 
  such that $a_n\in V_n$ and $f(a_n)\in U_n$. 
  This means that $a_n\to x_0$, $a_n\neq x_0$, $f(a_n)\to \ell$ 
  as desired.
  
  Conversely, if there exists such a sequence $a_n\to x_0$, then 
  for every $V$ neighborhood of $x_0$, it is eventually true that $a_n\in V$.
  And if $f(a_n)\to \ell$, for every $U$ neighborhood of $\ell$,
  it is also eventually true that $f(a_n)\in U$.
  Thus there exists $n$ such that $a_n\in V$ and $f(a_n)\in U$, confirming 
  that $f(x)\in U$ frequently. 
\end{proof}

Very often we will consider the limit points of a sequence $a_n$
as $n\to +\infty$.
In this case, the previous characterization tells us that $\ell$ 
is a limit point of $a_n$ as $n\to+\infty$ if there exists 
a sequence of indices $n_k\to+\infty$ such that $a_{n_k}\to \ell$.
We could also assume $n_k$ to be strictly increasing since if 
$n_k\to \infty$, we can extract from $n_k$ a strictly increasing subsequence. 
Thus the set of limit points of a sequence $a_n$ 
corresponds to the set of limits of all possible 
subsequences of $a_n$.

\begin{example}
  The sequence $a_n=(-1)^n$ has two limit points:
  \[
    \limsup_{n\to +\infty}(-1)^n = 1, \qquad \liminf_{n\to+\infty} (-1)^n = -1.
  \]
  Indeed, the subsequence of terms with even indices is constant $1$, while
  that of terms with odd indices is constant $-1$.
  There cannot be other limit points because if there were $a_{n_k}\to \ell$,
  then certainly $a_{n_k}=1$ frequently or $a_{n_k}=-1$
  frequently, from which either $\ell=1$ or $\ell=-1$.
\end{example}
%
\begin{example}
  Since $\#\QQ=\#\NN$, there exists a function $\vec a \colon \NN \to \QQ$
  that is surjective. Using the density property of rational numbers,
  it can be shown that the set of limit points of the corresponding sequence
  $a_n = \vec a(n)$ is the entire $\bar \RR$.
\end{example}

\begin{exercise}
  Consider the sequence $a_n = \sqrt n -\lfloor \sqrt n\rfloor$.
  Calculate 
  \[
    \liminf_{n\to+\infty} a_n, \qquad 
    \limsup_{n\to+\infty} a_n.
  \]
  What is the set of limit points of $a_n$ as $n\to +\infty$?
\end{exercise}

%
\begin{theorem}[properties of the upper/lower limit]
Let $f\colon A\subset \RR \to \RR$ and let $x_0$ be an accumulation point for
$A$.
Let $L$ be the set of limit points of $f(x)$ as $x\to x_0$. 
Then:
\begin{enumerate}
  \item %1
  $L\neq \emptyset$;

  \item %2
  $\displaystyle\limsup_{x\to x_0} f(x) \ge \liminf_{x\to x_0} f(x)$;

  \item %3
  if $\displaystyle\limsup_{x\to x_0} f(x) = \liminf_{x\to x_0} f(x) = \ell$, 
  then $\displaystyle\lim_{x\to x_0} f(x) = \ell$;

  \item %4
  the set of limit points is closed under taking limits:
  if $\ell_k\in L$ and $\ell_k \to \ell$ for some $\ell \in \bar \RR$, 
  then $\ell \in L$;

  \item %5
    $\displaystyle\limsup_{x\to x_0} f(x)$ and $\displaystyle\liminf_{x\to x_0}
  f(x)$ are limit points;

  \item %6
  the condition
  $\displaystyle\limsup_{x\to x_0} f(x) = \ell$ is equivalent to
  \[
  \begin{cases}
   \forall q > \ell \colon f(x) < q \text{ eventually,} \\
   \forall q < \ell \colon f(x) > q \text{ frequently,}
  \end{cases}
  \]
    and the condition $\displaystyle\liminf_{x\to x_0} f(x) = \ell$ is
  equivalent to
  \[
  \begin{cases}
  \forall q > \ell \colon f(x) < q \text{ frequently,} \\
  \forall q < \ell \colon f(x) > q \text{ eventually;}
  \end{cases}
  \]

  \item %7
  if $a_n$ is a sequence 
  \[
    \limsup_{n\to +\infty} a_n = \lim_{n\to +\infty} \sup_{k\ge n} a_k,
    \quad
    \liminf_{n\to +\infty} a_n = \lim_{n\to +\infty} \inf_{k\ge n} a_k.
  \]
 \end{enumerate}
\end{theorem}
%
\begin{proof}
  Regarding point 1.\
  the Bolzano-Weierstrass theorem~\ref{th:bolzano_weierstrass}
  guarantees that $L\neq \emptyset$, and
  therefore $\sup L \ge \inf L$, from which point 2 follows.

  For point 3.\ if $\limsup = \liminf = \ell$, it means that $L=\ENCLOSE{\ell}$
  has only one element. But by the corollary to the Bolzano-Weierstrass theorem,
  we know that from every sequence $f(a_n)$ with $a_n\to x_0$, 
  it is possible to extract a regular subsequence. 
  The limit of such a subsequence must be an element of $L$,
  and therefore it can only be $\ell$. 
  Then, by Proposition~\ref{prop:convergenza},
  we deduce that the entire sequence has limit $\ell$.

  For point 4.\ let $\ell_k\in L$ be a sequence of limit points
  and let $\ell\in \bar \RR$ be such that $\ell_k \to \ell$ as $k\to +\infty$.
  Now consider the sequence $U_n$ of neighborhoods of $\ell$ 
  given by Proposition~\ref{prop:base_numerabile}.
  Since $\ell_k\to \ell$, for every $n$, there must exist $k_n$ such 
  that $\ell_{k_n}\in U_n$. Without loss of generality, we can 
  directly assume that $\ell_n\in U_n$.
  Consider also the neighborhoods $V_n$ of $x_0$ given again 
  by Proposition~\ref{prop:base_numerabile}.
  For each $\ell_n \in L$, there must exist a sequence $a_k\to x_0$ 
  with $a_k\neq x_0$ such that $f(a_k)\to \ell_n$.
  But then there certainly exists $k=k_n$ such that $a_{k_n}\in V_n$ 
  and $f(a_{k_n})\in U_n$. 
  Setting $b_n = a_{k_n}$, we have found a sequence $b_n$ 
  such that $b_n\to x_0$ (since $b_n\in V_n$) and $f(b_n)\to \ell$
  since $f(b_n)\in U_n$. Thus $\ell$ is also a limit point. 

  For point 5,
  since $L\neq \emptyset$, the characterization of the $\sup$
    guarantees that there exists a sequence $\ell_n\in L$ such that $\ell_n \to
  \sup L$.
  But then, by the previous point, we must have $\sup L \in L$.
  The same holds for the $\inf$.

  For point 6.\ we demonstrate it for the $\limsup$ (for the $\liminf$,
  it will be analogous). If $\limsup f(x) = \ell$ and if it were frequently
  $f(x) \ge q > \ell$, then there would exist a sequence $a_n\to x_0$, 
  $a_n\neq x_0$ such that $f(a_n)\ge q$.
  This sequence would have a regular subsequence
  that has a limit $\ge q> \ell$, which is absurd.
  Thus it must be eventually true that $f(x) < q$ for every $q>\ell$.
  Conversely, if it were eventually true that $f(x) \le q < \ell$ for every
  sequence $a_n\to x_0$ with $a_n\neq x_0$, it would have to be 
  $f(a_n) \le q$, and this contradicts the fact that $\ell>q$ 
  is a limit point.

  Conversely, if for every $q>\ell$, it is eventually true that $f(x)<q$,
  it means that for every sequence $a_n\to x_0$, it is eventually true that
  $f(a_n) < q$, and therefore if $f(a_n)\to \ell$,
  it must be $\ell\le q$. 
  Thus every $q>\ell$ is an upper bound of $L$, and $\sup L\le \ell$.
  If then for every $q<\ell$, it is frequently true that $f(x)\ge q$,
  it means that there exists a sequence $a_n\to x_0$, $a_n\neq x_0$ 
  such that $f(a_n)\ge q$. 
  By the Bolzano-Weierstrass theorem, it will be possible to find 
  a convergent subsequence $f(a_{n_k})\to \ell'\ge q$.
  Thus for every $q<\ell$, it results
  $\sup L \ge q$, from which ultimately $\sup L = \ell$. 

  For point 7.\ we demonstrate it for the $\limsup$ (for the $\liminf$,
  it will be analogous).
  Setting $A_n = \sup_{k\ge n} a_k$, observe that $A_n$ is decreasing
  since as $n$ increases, the set
  $\ENCLOSE{k\colon k \ge n}$ decreases (in the sense of set inclusion),
  and therefore
  the $\sup$ does not increase.
  Thus $\ell=\lim A_n$ certainly exists in $\bar \RR$.
  Moreover, by the definition of limit, for every $\ell'>\ell$,
  it must be eventually true that
  $A_n < \ell'$, and therefore
  (since obviously $a_n \le A_n$), it will also be
  $a_n < \ell'$ eventually.
  Conversely, given $\ell'<\ell$, it will be eventually true that
  $A_n > \ell'$. But this means that there exists $k>n$
  for which $a_k > \ell'$, and therefore
  it results $a_n > \ell'$ frequently.
  By the previous point, we can then conclude that
  $\ell = \limsup a_n$.
\end{proof}

\begin{theorem}[operations with $\limsup$ and $\liminf$]
    Let $f,g\colon A\subset \RR\to\RR$ and let $x_0$ be an accumulation point of
  $A$.
  \begin{enumerate}
    \item 
      As $x\to x_0$, we have 
      \begin{align*}
      \limsup_{x\to x_0} (-f(x)) &= -\liminf_{x\to x_0} f(x)\\ 
      \liminf_{x\to x_0} (-f(x)) &= -\limsup_{x\to x_0} f(x);
      \end{align*}
    \item
      if $\lambda \ge 0$, then
        \begin{align*}
        \limsup_{x\to x_0} (\lambda \cdot f(x)) &= \lambda \cdot \limsup_{x\to x_0} f(x), \\
        \liminf_{x\to x_0} (\lambda \cdot f(x)) &= \lambda \cdot \liminf_{x\to x_0} f(x);
        \end{align*}
    \item
    moreover
    \begin{gather*}
    \limsup_{x\to x_0} (f(x) + g(x)) \le \limsup_{x\to x_0} f(x) + \limsup_{x\to x_0} g(x),\\
    \liminf_{x\to x_0} (f(x) + g(x)) \ge \liminf_{x\to x_0} f(x) + \liminf_{x\to x_0} g(x);
    \end{gather*}
    \item
    and finally, if $f(x) \ge g(x)$, then
    \[
     \limsup_{x\to x_0} f(x) \ge \limsup_{x\to x_0} g(x), \qquad \liminf_{x\to x_0} f(x) \ge \liminf_{x\to x_0} g(x).
    \]
  \end{enumerate}
\end{theorem}
%
\begin{proof}
  For point 1.\ it suffices to observe that if the function $f(x)$ has $L$ as
  the set of limit points, then the function $-f(x)$ has $-L$ as limit
  points. Therefore $\sup (-L) = -\inf L$ and $\inf(-L) = -\sup L$.

  For point 2.\ note that if $L$ is the set of limit points
  of $f(x)$, then the set of limit points of $\lambda f(x)$ is $\lambda L$.
  If $\lambda \ge 0$, then $\sup \lambda L = \lambda \sup L$
    and $\inf \lambda L = \lambda \inf L$ (if $\lambda<0$, instead, $\inf$ and
  $\sup$
  are exchanged, as in point 1).

    For point 3.\ consider the case of the $\limsup$ (for the $\liminf$, it will
  be analogous).
  If $\ell = \limsup(f(x)+g(x))$, it means that there exists a sequence
  $a_n\to x_0$ such that $f(a_n)+g(a_n)\to \ell$. 
  But I can extract a subsequence
  such that the first addend $f(a_{n_k})$ also has a limit. 
  And then I can extract
  a sub-subsequence so that the second addend $g(a_{n_{k_j}})$
  also has a limit. 
  Thus we will have $f(a_{n_{k_j}}) \to \ell_1$, $g(a_{n_{k_j}} \to \ell_2$
  with $\ell_1+\ell_2=\ell$.
    But then by the definition of $\limsup$, we will have $\limsup f(x) \ge
  \ell_1$
    and $\limsup g(x) \ge \ell_2$, from which $\limsup f(x) + \limsup g(x) \ge
  \ell_1+\ell_2 = \ell$.

  For point 4.\ note that if $\ell = \limsup f(x)$, then
  for every $q>\ell$, it is eventually true that
  $f(x)<q$, and consequently also $g(x) < q$ eventually.
  Thus $\limsup g(x) \le \ell$.
  If instead we set $\ell = \liminf h(x)$, then
  for every $q<l$, it is eventually true that $g(x) \ge q$, but
  then also $f(x)\ge q$ eventually, and therefore
  $\liminf f(x) \ge \ell$.
\end{proof}
    
\begin{theorem}[Cesàro convergence]
  \label{th:criterio_cesaro}%
  \mymark{*}%
  \mymargin{Cesàro}%
\index{Cesàro}%
  \index{criterion!of Cesàro ratio}%
  \index{sum!of Cesàro}%
  Let $a_n$ be a sequence.
  \begin{enumerate}
  \item
    If
    $   a_n \to \ell \in \bar \RR$,
    then
    \[
    \frac 1 n \cdot \displaystyle\sum_{k=1}^n a_k \to \ell.
    \]
  
  \item
    If $a_n>0$,
    $\displaystyle\frac{a_{n+1}}{a_n} \to \ell \in [0,+\infty]$,
    then
    $\displaystyle \sqrt[n]{a_n}\to \ell$.
  \end{enumerate}
  \end{theorem}
  %
  \begin{proof}
  Let us prove the first point.
  For every $q<\ell$, since $a_n\to \ell$, there must exist $N$ such that 
  $a_n\ge q$ if $n> N$. 
  Thus 
  \[
    b_n = \frac 1 n \sum_{k=1}^{n} a_k 
    \ge \frac 1 n \sum_{k=1}^N a_k + \frac 1 n \sum_{k=N+1}^n q 
    = \frac{\sum_{k=1}^N a_k}{n} + \frac{n-N}{n}\cdot q \to q.
  \]
  This means that for every $q<\ell$, $\liminf b_n \ge q$, and 
  thus $\liminf b_n \ge \ell$.
  Conversely, for every $q>\ell$, reasoning analogously, we obtain 
  that $b_n$ does not exceed a sequence that tends to $q$.
  Thus $\limsup b_n\le \ell$. 
  Putting things together, we deduce that $\lim b_n = \ell$.
    
  Point 2 can be reduced to point 1.
  Indeed,
  \[
    \ln \sqrt[n]{a_n}
    = \ln \sqrt[n]{a_0\cdot \frac{a_1}{a_0} \cdots \frac{a_n}{a_{n-1}}}
    = \frac{\ln a_0 + \ln \frac{a_1}{a_0} + \dots + \ln \frac{a_n}{a_{n-1}}}{n}.
  \]
  Setting $x_n = \ln \frac{a_n}{a_{n-1}}$,
  if $\frac{a_n}{a_{n-1}}\to \ell$,
    we have $x_n\to \ln \ell$ (interpreting $\ln 0 = -\infty$ and $\ln
  (+\infty)=+\infty$), and thus
  \[
    \ln \sqrt[n]{a_n} = \frac{\ln a_0}{n} + \frac{x_1 + \dots + x_n}{n}
    \to \ln \ell.
  \]
  Taking the exponential of both sides, we find the desired result.
  \end{proof}
  
  \begin{exercise}\label{ex:7340098}
  Apply the previous result to
  verify that
  \[
     \lim \sqrt[n]{n} = 1
  \]
 and
 \index{factorial!asymptotic estimate}%
 \[
   \lim \frac{n}{\sqrt[n]{n!}} = e.
 \]
 \end{exercise}
 \mynote{In exercise~\ref{ex:7340098}, it is shown that 
 \[
  \sqrt[n]{n!} \sim \frac{n}{e}  
  \qquad\text{as $n\to +\infty$.}
 \]
Note that it is not possible to raise both sides of this asymptotic estimate to
the $n$th power.
Indeed, in Theorem~\ref{th:stirling}, we will find this asymptotic estimate for
the factorial:
 \[
  n! \sim \sqrt{2\pi n} \cdot \frac{n^n}{e^n}.
 \]
 See example~\ref{ex:498124} for a similar result.%
 }

 \begin{exercise}
 Consider the sequence
 \[
 a_n =
 \begin{cases}
    2^n &\text{if $n$ is even},\\
    n\cdot 2^n &\text{if $n$ is odd}.
 \end{cases}
 \]
 Verify that the sequence $a_n$
  can be tested with the root criterion but
  not with the ratio criterion.
  Use this example to show that the implications
  stated in Theorem~\ref{th:criterio_cesaro} cannot
  be reversed.
\end{exercise}  

\subsection{the zero theorem}

\begin{theorem}[zero theorem]
\mymark{***}%
\mymargin{zero theorem}%
\index{theorem!zero theorem}%
\label{th:zeri}%
Let $f\colon[a,b] \to \RR$ be a continuous function
such that $f(a)\le 0$ and $f(b)\ge 0$.
\mynote{If $b<a$, it is understood that $[a,b]=[b,a]$.}
Then there exists
 $c\in [a,b]$ such that $f(c)=0$.
\end{theorem}

\begin{proof}
\mymark{***}
The proof we adopt is of particular importance as
it not only allows us to prove the existence of the point $c$ that solves
$f(x)=0$
but also presents us with
an algorithm, the \emph{bisection method}%
\mymargin{bisection method}%
\index{method!bisection method},
\index{bisection!method of}
that can be effectively used to approximate
such a solution.

We can assume without loss of generality that $a<b$.
Set $A_0 = a$, $B_0= b$, and consider the midpoint $C_0 = (A_0+B_0)/2$.
Choose between the two intervals $[A_0, C_0]$ and $[C_0,B_0]$ the one for which
the sign at the two ends is opposite (or, in a fortunate case, zero).
More precisely, if $f(C_0)\ge 0$, set $[A_1,B_1] = [A_0,C_0]$; otherwise,
choose $[A_1,B_1] = [C_0,B_0]$ so that, in any case,
$f(A_1)\le 0$, $f(B_1)\ge 0$.

Consider the midpoint $C_1$ of the new interval $[A_1,B_1]$ and repeat
the process indefinitely. What we obtain are two sequences
$A_n$, $B_n$ with these properties (which could be proven by induction):
\begin{enumerate}
\item $A_n < B_n$, $B_n - A_n = (b-a)/2^n$;
\item $A_n$ is increasing, $B_n$ is decreasing;
\item $f(A_n)\le 0$, $f(B_n)\ge 0$.
\end{enumerate}

Since $A_n$ is monotonic, we know that $A_n$ converges $A_n\to c$.
Moreover, since $A_n \in [a,b]$, also $c\in [a,b]$ (by the permanence of
the sign of the sequences $A_n-a$ and $b-A_n$).
Passing to the limit in the equality $B_n = A_n + (b-a)/2^n$,
we obtain that also $B_n \to c$. Since $f$ is continuous,
we will have
\[
f(A_n) \to f(c), \qquad
f(B_n) \to f(c).
\]
But $f(A_n)\le 0$, and therefore by the permanence of the sign, also $f(c)\le
0$.
On the other hand, $f(B_n) \ge 0$, and therefore $f(c)\ge 0$.
Thus we obtain $f(c) = 0$, as we wanted to prove.
\end{proof}

\begin{example}\label{ex:75445}
Solve the equation
\[
  x^5-x-1=0.
\]
\end{example}
%
\begin{proof}[Solution]
Let $f(x) = x^5-x-1$. It is clear that the function $f\colon \RR\to \RR$
is continuous (as it is a composition of continuous functions).
Observe that $f(0) = -1$ and $f(2)=29$, so the function
satisfies the hypotheses of the zero theorem on the interval $[0,2]$.
We know, therefore, that the equation in question has at least one solution
in this interval.

Using the bisection method, we can determine a solution
with arbitrary precision. Setting $A_0=0$, $B_0=2$, we have verified that
$f(A_0)<0$ and $f(B_0)>0$.
Take the midpoint
$C_0=1$ and calculate the function: $f(C_0)=-1 < 0$. We know
then that a solution must be contained in the interval
$[A_1,B_1] = [C_0,B_0] = [1,2]$ because the hypotheses
of the zero theorem also hold in this interval.
The midpoint of this interval is $C_1=3/2 = 1.5$,
and it turns out that $f(3/2) = 163/32>0$, so the next interval
we will consider is $[A_2,B_2]=[1,3/2]$.
To avoid working with too many decimal places, instead of dividing
this interval exactly in half, we consider an intermediate point
$C_2 = 6/5 = 1.2$ where $f(C_2)=901/3125>0$.
We know then that a solution is contained in the interval
$[A_3,B_3] = [1,1.2]$. Take the midpoint $C_3=11/10=1.1$
and find $f(C_3) = -48949/10^5 <0$. We have thus obtained
that there exists $x\in (1.1,1.2)$ such that $f(x)=0$. We know then
that $\abs{x-1.15} < 0.05$, i.e., we have found $x$ with an error
less than $0.05$.

With much patience, one can proceed
with the bisection method until verifying
that $f(116/10^2)$ $=$ $-596583424/10^{10}<0$ and
$f(117/10^2)=224480357/10^{10}>0$, from which
it is obtained that a solution is contained between $1.16$ and $1.17$ with an
error
less than $0.005$.
With a calculator (see, for example, the code on page \pageref{code:bisection}),
more significant digits can be obtained: $x=1.1673039782614187\ldots$
\end{proof}

\begin{table}
\begin{center}
\begin{tabular}{r}
$\sqrt 2 \approx $ \ttfamily\footnotesize 
1.4142135623 7309504880 1688724209 6980785696 7187537694 \\
\ttfamily\footnotesize
8073176679 7379907324 7846210703 8850387534 3276415727 \\
% 3501384623 0912297024 9248360558 5073721264 4121497099 \\ \ttfamily\footnotesize
% 9358314132 2266592750 5592755799 9505011527 8206057147 \\ \ttfamily\footnotesize
% 0109559971 6059702745 3459686201 4728517418 6408891986 \\ \ttfamily\footnotesize
% 0955232923 0484308714 3214508397 6260362799 5251407989 \\ \ttfamily\footnotesize
% 6872533965 4633180882 9640620615 2583523950 5474575028 \\ \ttfamily\footnotesize
% 7759961729 8355752203 3753185701 1354374603 4084988471 \\ \ttfamily\footnotesize
% 6038689997 0699004815 0305440277 9031645424 7823068492 \\ \ttfamily\footnotesize
% 9369186215 8057846311 1596668713 0130156185 6898723723 \\ \ttfamily\footnotesize
% 5288509264 8612494977 1542183342 0428568606 0146824720 \\ \ttfamily\footnotesize
% 7714358548 7415565706 9677653720 2264854470 1585880162 \\ \ttfamily\footnotesize
% 0758474922 6572260020 8558446652 1458398893 9443709265 \\ \ttfamily\footnotesize
% 9180031138 8246468157 0826301005 9485870400 3186480342 \\ \ttfamily\footnotesize
% 1948972782 9064104507 2636881313 7398552561 1732204024 \\ \ttfamily\footnotesize
% 5091227700 2269411275 7362728049 5738108967 5040183698 \\ \ttfamily\footnotesize
% 6836845072 5799364729 0607629969 4138047565 4823728997 \\ \ttfamily\footnotesize
% 1803268024 7442062926 9124859052 1810044598 4215059112 \\ \ttfamily\footnotesize
% 0249441341 7285314781 0580360337 1077309182 8693147101 \\ \ttfamily\footnotesize
% 7111168391 6581726889 4197587165 8215212822 9518488472 \\
\\
$\sqrt 3 \approx$ \ttfamily\footnotesize 
1.7320508075 6887729352 7446341505 8723669428 0525381038 \\
\ttfamily\footnotesize
0628055806 9794519330 1690880003 7081146186 7572485757 \\
\\
$\phi = \frac{\sqrt 5+1}{2} \approx $ \ttfamily\footnotesize 
1.6180339887 4989484820 4586834365 6381177203 0917980576 \\
\ttfamily\footnotesize
2862135448 6227052604 6281890244 9707207204 1893911375
  \end{tabular}
\end{center}
\caption{The first 100 decimal digits of some
constants calculated with the bisection method used in the proof
of Theorem~\ref{th:zeri}.
See the code on page~\pageref{code:bisection}.}
\label{fig:cifre_sqrt2}
\index{$\sqrt 2$!decimal digits}
\index{digits!$\sqrt 2$}
\end{table}

\begin{theorem}[intermediate value theorem]
\label{th:valori_intermedi}%
\mymark{**}%
\mymargin{intermediate value property}%
\index{property!intermediate value}%
\index{theorem!intermediate value}%
Let $I\subset \RR$ be an interval and $f\colon I \to \RR$ a
continuous function.
Then if $f$ assumes two values $y_1$ and $y_2$, it also assumes all intermediate
values between $y_1$ and $y_2$.
In other words: a continuous function
maps intervals into intervals.
\end{theorem}
%
\begin{proof}
If $y_1$ and $y_2$ are values assumed by $f$, it means
that there exist $x_1,x_2 \in I$ such that $f(x_1)= y_1$ and $f(x_2)=y_2$.
Then, given $y$, consider the function $g(x) = f(x)-y$.
If $y$ is intermediate between $y_1$ and $y_2$, the function $g$ will
assume opposite signs at $x_1$ and $x_2$, and thus, by the zero theorem,
there must be a point $x$ where $g$ vanishes. At this point,
we have $f(x)=y$, as we wanted to prove.
\end{proof}

\begin{lemma}
\label{ex:inversa_monotona}%
If $I$ is an interval of $\RR$, every function $f\colon I \to \RR$
that is injective and continuous is strictly monotonic.
\end{lemma}
%
\begin{proof}
It can be observed that a function is strictly monotonic if it maintains
intermediate values, i.e., if given three points
$x<y<z$, it always results that $f(y)$ is an intermediate value between $f(x)$
and $f(z)$:
\[
  f(x)< f(y) <f(z) \qquad\text{or} \qquad f(x)> f(y) > f(z).
\]
If this did not happen, for example, if it were $f(y)>f(z)>f(x)$ with $x<y<z$,
then by the continuity of $f$, there would have to exist an intermediate value
between $x$ and $y$ where the function assumes the value $f(z)$. But then the
function
would not be injective.
\end{proof}

%%%%%%%%%%%%%%%%%%%%%%%%
%%%%%%%%%%%%%%%%%%%%%%%%
%%%%%%%%%%%%%%%%%%%%%%%%
%%%%%%%%%%%%%%%%%%%%%%%%